\subsection{معرفی و انگیزه}
این ماژول در چارچوب پژوهش حاضر پیاده‌سازی و ارزیابی شده و فاقد مقاله مستقل جداگانه در ادبیات تحقیق است.

در بسیاری از سناریوهای فشرده‌سازی لایه‌های ترنسفورمر، تحمیل یک رتبه ثابت از پیش تعیین‌شده به همه لایه‌ها غیرواقع‌بینانه است؛ چرا که لایه‌های مختلف دارای زوال طیفی متفاوتی در مقادیر منفرد خود هستند. برخی لایه‌ها با نگه‌داشتن تعداد اندکی از مقادیر منفرد بیش از ۹۵ درصد از انرژی وزن را حفظ می‌کنند، در حالی که لایه‌های دیگر نیاز به رتبه بالاتری برای حفظ کارکرد خود دارند. در نتیجه، تنظیم مبتنی بر رتبه اغلب یا به فشرده‌سازی ناکافی یا به تخریب غیرقابل جبران دقت می‌انجامد.

ماژول فشرده‌ساز \lr{ERC-SVD} با نام ثبتی \lr{erc} و \lr{erc\_svd} («تجزیه مقادیر منفرد با خطای کنترل‌شده» یا \lr{Error-Controlled SVD})، با اتخاذ رویکرد دوگان، هدف فشرده‌سازی را بر پایه یک سقف خطای نسبی فروبنیوس مجاز (\lr{Bounded Frobenius Truncation Error}) تعریف می‌کند. در این راهبرد، به ازای هر لایه، کمترین رتبه ممکن $r$ به گونه‌ای انتخاب می‌شود که تضمین کند خطای برش مقادیر منفرد از سقف مشخص‌شده فراتر نرود. با این کار، توازن بهینه‌ای میان نرخ تراکم و پایداری عملکرد مدل برقرار می‌گردد.
